\documentclass[10pt,a4paper]{amsart}
\usepackage{amsmath,amssymb,amsthm,graphicx}
\usepackage{verbatim}
\newenvironment{reference}{\comment}{\endcomment}
%\newenvironment{reference}{}{}
\newenvironment{slogan}{\comment}{\endcomment}
\newenvironment{history}{\comment}{\endcomment}
% For commutative diagrams we use Xy-pic
\usepackage[all]{xy}

% We use 2cell for 2-commutative diagrams.
\xyoption{2cell}
% Theorem environments.
%
\theoremstyle{plain}
\newtheorem{theorem}[subsection]{Theorem}
\newtheorem{proposition}[subsection]{Proposition}
\newtheorem{lemma}[subsection]{Lemma}

\theoremstyle{definition}
\newtheorem{definition}[subsection]{Definition}
\newtheorem{example}[subsection]{Example}
\newtheorem{exercise}[subsection]{Exercise}
\newtheorem{situation}[subsection]{Situation}

\theoremstyle{remark}
\newtheorem{remark}[subsection]{Remark}
\newtheorem{remarks}[subsection]{Remarks}

\numberwithin{equation}{subsection}

% Macros
%
\def\lim{\mathop{\mathrm{lim}}\nolimits}
\def\colim{\mathop{\mathrm{colim}}\nolimits}
\def\Spec{\mathop{\mathrm{Spec}}}
\def\Hom{\mathop{\mathrm{Hom}}\nolimits}
\def\Ext{\mathop{\mathrm{Ext}}\nolimits}
\def\SheafHom{\mathop{\mathcal{H}\!\mathit{om}}\nolimits}
\def\SheafExt{\mathop{\mathcal{E}\!\mathit{xt}}\nolimits}
\def\Sch{\mathit{Sch}}
\def\Mor{\mathop{\mathrm{Mor}}\nolimits}
\def\Ob{\mathop{\mathrm{Ob}}\nolimits}
\def\Sh{\mathop{\mathit{Sh}}\nolimits}
\def\NL{\mathop{N\!L}\nolimits}
\def\CH{\mathop{\mathrm{CH}}\nolimits}
\def\proetale{{pro\text{-}\acute{e}tale}}
\def\etale{{\acute{e}tale}}
\def\QCoh{\mathit{QCoh}}
\def\Ker{\mathop{\mathrm{Ker}}}
\def\Im{\mathop{\mathrm{Im}}}
\def\Coker{\mathop{\mathrm{Coker}}}
\def\Coim{\mathop{\mathrm{Coim}}}

% Boxtimes
%
\DeclareMathSymbol{\boxtimes}{\mathbin}{AMSa}{"02}

%
% Macros for moduli stacks/spaces
%
\def\QCohstack{\mathcal{QC}\!\mathit{oh}}
\def\Cohstack{\mathcal{C}\!\mathit{oh}}
\def\Spacesstack{\mathcal{S}\!\mathit{paces}}
\def\Quotfunctor{\mathrm{Quot}}
\def\Hilbfunctor{\mathrm{Hilb}}
\def\Curvesstack{\mathcal{C}\!\mathit{urves}}
\def\Polarizedstack{\mathcal{P}\!\mathit{olarized}}
\def\Complexesstack{\mathcal{C}\!\mathit{omplexes}}
% \Pic is the operator that assigns to X its picard group, usage \Pic(X)
% \Picardstack_{X/B} denotes the Picard stack of X over B
% \Picardfunctor_{X/B} denotes the Picard functor of X over B
\def\Pic{\mathop{\mathrm{Pic}}\nolimits}
\def\Picardstack{\mathcal{P}\!\mathit{ic}}
\def\Picardfunctor{\mathrm{Pic}}
\def\Deformationcategory{\mathcal{D}\!\mathit{ef}}
\usepackage{cleveref}
\crefname{originallemma}{original lemma}{original lemmas}
\Crefname{originallemma}{Original lemma}{Original lemmas}
\makeatletter
\newlabel{lemma-affine-morphism-projection-ideal}{{0GF8}{source}{}{sourcepointer-1}{}}
\newlabel{lemma-affine-morphism-projection-ideal@cref}{{[originallemma][0][]0GF8}{[1][0][]source}{}{}{}}
\makeatother
\crefname{originallemma}{original lemma}{original lemmas}
\Crefname{originallemma}{Original lemma}{Original lemmas}
\makeatletter
\newlabel{lemma-affine-vanishing-higher-direct-images}{{073H}{source}{}{sourcepointer-2}{}}
\newlabel{lemma-affine-vanishing-higher-direct-images@cref}{{[originallemma][0][]073H}{[1][0][]source}{}{}{}}
\makeatother
\crefname{originallemma}{original lemma}{original lemmas}
\Crefname{originallemma}{Original lemma}{Original lemmas}
\makeatletter
\newlabel{lemma-colimits}{{073E}{source}{}{sourcepointer-3}{}}
\newlabel{lemma-colimits@cref}{{[originallemma][0][]073E}{[1][0][]source}{}{}{}}
\makeatother
\crefname{originallemma}{original lemma}{original lemmas}
\Crefname{originallemma}{Original lemma}{Original lemmas}
\makeatletter
\newlabel{lemma-directed-colimit-coherent}{{07UV}{source}{}{sourcepointer-4}{}}
\newlabel{lemma-directed-colimit-coherent@cref}{{[originallemma][0][]07UV}{[1][0][]source}{}{}{}}
\makeatother
\crefname{originallemma}{original lemma}{original lemmas}
\Crefname{originallemma}{Original lemma}{Original lemmas}
\makeatletter
\newlabel{lemma-finite-morphism-Noetherian}{{0GF7}{source}{}{sourcepointer-5}{}}
\newlabel{lemma-finite-morphism-Noetherian@cref}{{[originallemma][0][]0GF7}{[1][0][]source}{}{}{}}
\makeatother
\crefname{originallemma}{original lemma}{original lemmas}
\Crefname{originallemma}{Original lemma}{Original lemmas}
\makeatletter
\newlabel{lemma-property-higher-rank-cohomological}{{07UT}{source}{}{sourcepointer-6}{}}
\newlabel{lemma-property-higher-rank-cohomological@cref}{{[originallemma][0][]07UT}{[1][0][]source}{}{}{}}
\makeatother
\crefname{originalproposition}{original proposition}{original propositions}
\Crefname{originalproposition}{Original proposition}{Original propositions}
\makeatletter
\newlabel{proposition-vanishing-affine}{{07V6}{source}{}{sourcepointer-7}{}}
\newlabel{proposition-vanishing-affine@cref}{{[originalproposition][0][]07V6}{[1][0][]source}{}{}{}}
\makeatother
\crefname{originallemma}{original lemma}{original lemmas}
\Crefname{originallemma}{Original lemma}{Original lemmas}
\makeatletter
\newlabel{sites-cohomology-lemma-apply-Leray}{{0733}{source}{}{sourcepointer-8}{}}
\newlabel{sites-cohomology-lemma-apply-Leray@cref}{{[originallemma][0][]0733}{[1][0][]source}{}{}{}}
\makeatother
\title{A Noetherian algebraic space receiving a finite surjective morphism from an affine space is affine}
\author{The Stacks Project Authors}
\date{stacks-project:spaces-cohomology.tex:local-snapshot-20261012; GNU FDL 1.2 or later}
\begin{document}
\maketitle
\noindent\textit{Editorial note.} All delivered mathematical text is the author's own, copied from the pinned Stacks Project chapter file spaces-cohomology.tex; no statement or proof step was rewritten, reordered or removed.\par\medskip
\noindent\textit{Editorial note.} Preamble provenance: the pinned chapter file opens with \textbackslash{}input\{preamble\} and carries no preamble text of its own, so the project's single Stacks-wide preamble staged at sources/stacks-spaces-cohomology/preamble.tex is excerpted verbatim as three spans, lines 8-12 (verbatim package and the project's own reference/slogan/history display environments), lines 14-18 (Xy-pic setup) and lines 154-217 (theorem declarations and mathematical macros). The wrapper class and the amsmath/amssymb/amsthm/graphicx base packages are supplied by the assembler. No macro was invented or redefined.\par\medskip
\noindent\textit{Editorial note.} Same-chapter pointers are located at their \textbackslash{}label lines inside the pinned chapter file; cross-chapter pointers are located in the project's own tag list staged in the snapshot (sources/stacks-chow/tags.txt), at the verbatim lines carrying their canonical tags.\par\medskip
\noindent\textit{Allowed standard prerequisites (editorial).}\begin{itemize}
\item Extension from coherent to quasi-coherent modules by filtered colimits, and the affineness criterion by vanishing of \$H\textasciicircum{}1\$. Same chapter spaces-cohomology.tex, labels lemma-directed-colimit-coherent and lemma-colimits (invoked at proof line 3434) and label proposition-vanishing-affine (invoked at proof line 3436). They reduce affineness of \$X\$ to the vanishing of \$H\textasciicircum{}1(X, \textbackslash{}mathcal\{F\}) = 0\$ for all coherent \$\textbackslash{}mathcal\{F\}\$, which is what the rest of the proof establishes.
\item The criterion reducing a property of coherent sheaves on a Noetherian space to irreducible supports. Same chapter spaces-cohomology.tex, label lemma-property-higher-rank-cohomological (invoked at proof lines 3444 and 3453). Conditions (1) and (2) are verified directly for the property \$\textbackslash{}mathcal\{P\}(\textbackslash{}mathcal\{F\}) \textbackslash{}Leftrightarrow H\textasciicircum{}1(X, \textbackslash{}mathcal\{F\}) = 0\$, and condition (3) is met by the sheaf \$\textbackslash{}mathcal\{G\} = f'\_*\textbackslash{}mathcal\{O\}\_\{Z'\}\$.
\item The auxiliary finite morphism over a reduced irreducible closed subspace, and the projection formula for an affine morphism. Same chapter spaces-cohomology.tex, labels lemma-finite-morphism-Noetherian (invoked at proof lines 3450 and 3454) and lemma-affine-morphism-projection-ideal (invoked at proof line 3460). The first supplies \$Z'\$ with \$i' : Z' \textbackslash{}to Y\$, \$f' : Z' \textbackslash{}to Z\$ and the support property (3)(a); the second identifies \$\textbackslash{}mathcal\{G\}' = \textbackslash{}mathcal\{I\}\textbackslash{}mathcal\{G\}\$ with \$f'\_*\textbackslash{}mathcal\{I\}'\$.
\item Vanishing of the higher direct images of the affine morphisms \$f'\$ and \$i\$, with Leray. Same chapter spaces-cohomology.tex, label lemma-affine-vanishing-higher-direct-images (invoked at proof lines 3470 and 3477) and Cohomology on Sites, label sites-cohomology-lemma-apply-Leray (invoked at proof line 3474). Since \$f'\$ and \$i\$ are affine, their higher direct images vanish; together with the vanishing of \$H\textasciicircum{}1(Z', \textbackslash{}mathcal\{I\}')\$ on the affine space \$Z'\$ and Leray, this gives \$H\textasciicircum{}1(X, i\_*\textbackslash{}mathcal\{G\}') = 0\$, which is condition (3)(b).
\end{itemize}\medskip
\ifdefined\hypertarget\hypertarget{sourcepointer-1}{}\fi
\noindent\textit{Source reference.} Original lemma 0GF8: sources/stacks-spaces-cohomology/pinned-original-source.tex, source lines 3386--3386. Navigation pointer for \textbackslash{}ref\{lemma-affine-morphism-projection-ideal\} used inside the counted span (line 3460); the argument that consumes it is retained in the span.\par\medskip
\ifdefined\hypertarget\hypertarget{sourcepointer-2}{}\fi
\noindent\textit{Source reference.} Original lemma 073H: sources/stacks-spaces-cohomology/pinned-original-source.tex, source lines 1522--1522. Navigation pointer for \textbackslash{}ref\{lemma-affine-vanishing-higher-direct-images\} used inside the counted span (line 3470); the argument that consumes it is retained in the span.\par\medskip
\ifdefined\hypertarget\hypertarget{sourcepointer-3}{}\fi
\noindent\textit{Source reference.} Original lemma 073E: sources/stacks-spaces-cohomology/pinned-original-source.tex, source lines 667--667. Navigation pointer for \textbackslash{}ref\{lemma-colimits\} used inside the counted span (line 3434); the argument that consumes it is retained in the span.\par\medskip
\ifdefined\hypertarget\hypertarget{sourcepointer-4}{}\fi
\noindent\textit{Source reference.} Original lemma 07UV: sources/stacks-spaces-cohomology/pinned-original-source.tex, source lines 2866--2866. Navigation pointer for \textbackslash{}ref\{lemma-directed-colimit-coherent\} used inside the counted span (line 3434); the argument that consumes it is retained in the span.\par\medskip
\ifdefined\hypertarget\hypertarget{sourcepointer-5}{}\fi
\noindent\textit{Source reference.} Original lemma 0GF7: sources/stacks-spaces-cohomology/pinned-original-source.tex, source lines 3351--3351. Navigation pointer for \textbackslash{}ref\{lemma-finite-morphism-Noetherian\} used inside the counted span (line 3450); the argument that consumes it is retained in the span.\par\medskip
\ifdefined\hypertarget\hypertarget{sourcepointer-6}{}\fi
\noindent\textit{Source reference.} Original lemma 07UT: sources/stacks-spaces-cohomology/pinned-original-source.tex, source lines 2652--2652. Navigation pointer for \textbackslash{}ref\{lemma-property-higher-rank-cohomological\} used inside the counted span (line 3444); the argument that consumes it is retained in the span.\par\medskip
\ifdefined\hypertarget\hypertarget{sourcepointer-7}{}\fi
\noindent\textit{Source reference.} Original proposition 07V6: sources/stacks-spaces-cohomology/pinned-original-source.tex, source lines 3160--3160. Navigation pointer for \textbackslash{}ref\{proposition-vanishing-affine\} used inside the counted span (line 3436); the argument that consumes it is retained in the span.\par\medskip
\ifdefined\hypertarget\hypertarget{sourcepointer-8}{}\fi
\noindent\textit{Source reference.} Original lemma 0733: sources/stacks-chow/tags.txt, source lines 9047--9047. Navigation pointer for \textbackslash{}ref\{sites-cohomology-lemma-apply-Leray\} used inside the counted span (line 3474); the argument that consumes it is retained in the span.\par\medskip
\begin{lemma}
\label{lemma-image-affine-finite-morphism-affine-Noetherian}
Let $S$ be a scheme. Let $f : Y \to X$ be a morphism of algebraic
spaces over $S$. Assume
\begin{enumerate}
\item $f$ finite,
\item $f$ surjective,
\item $Y$ affine, and
\item $X$ Noetherian.
\end{enumerate}
Then $X$ is affine.
\end{lemma}

\begin{proof}
We will prove that under the assumptions of the lemma for any coherent
$\mathcal{O}_X$-module $\mathcal{F}$ we have $H^1(X, \mathcal{F}) = 0$.
This implies that $H^1(X, \mathcal{F}) = 0$ for every quasi-coherent
$\mathcal{O}_X$-module $\mathcal{F}$ by
Lemmas \ref{lemma-directed-colimit-coherent} and \ref{lemma-colimits}.
Then it follows that $X$ is affine from
Proposition \ref{proposition-vanishing-affine}.

\medskip\noindent
Let $\mathcal{P}$ be the property of coherent sheaves
$\mathcal{F}$ on $X$ defined by the rule
$$
\mathcal{P}(\mathcal{F}) \Leftrightarrow H^1(X, \mathcal{F}) = 0.
$$
We are going to apply Lemma \ref{lemma-property-higher-rank-cohomological}.
Thus we have to verify (1), (2) and (3) of that lemma for $\mathcal{P}$.
Property (1) follows from the long exact cohomology sequence associated
to a short exact sequence of sheaves. Property (2) follows since
$H^1(X, -)$ is an additive functor. To see (3) let $i : Z \to X$ be
a reduced closed subspace with $|Z|$ irreducible. Let $i' : Z' \to Y$
and $f' : Z' \to Z$ be as in Lemma \ref{lemma-finite-morphism-Noetherian}
and set $\mathcal{G} = f'_*\mathcal{O}_{Z'}$. We claim that
$\mathcal{G}$ satisfies properties (3)(a) and (3)(b) of
Lemma \ref{lemma-property-higher-rank-cohomological}
which will finish the proof. Property (3)(a) we have seen in
Lemma \ref{lemma-finite-morphism-Noetherian}. To see (3)(b) let
$\mathcal{I}$ be a nonzero quasi-coherent sheaf of ideals on $Z$.
Denote $\mathcal{I}' \subset \mathcal{O}_{Z'}$ the quasi-coherent
ideal $(f')^{-1}\mathcal{I} \mathcal{O}_{Z'}$, i.e., the
image of $(f')^*\mathcal{I} \to \mathcal{O}_{Z'}$.
By Lemma \ref{lemma-affine-morphism-projection-ideal} we have
$f_*\mathcal{I}' = \mathcal{I} \mathcal{G}$.
We claim the common value
$\mathcal{G}' = \mathcal{I} \mathcal{G} = f'_*\mathcal{I}'$
satisfies the condition expressed in (3)(b).
First, it is clear that the support of $\mathcal{G}/\mathcal{G}'$
is contained in the support of $\mathcal{O}_Z/\mathcal{I}$
which is a proper subspace of $|Z|$ as $\mathcal{I}$ is a
nonzero ideal sheaf on the reduced and irreducible algebraic space $Z$.
The morphism $f'$ is affine, hence $R^1f'_*\mathcal{I}' = 0$ by
Lemma \ref{lemma-affine-vanishing-higher-direct-images}.
As $Z'$ is affine (as a closed subscheme of an affine scheme)
we have $H^1(Z', \mathcal{I}') = 0$. Hence the Leray
spectral sequence (in the form
Cohomology on Sites, Lemma \ref{sites-cohomology-lemma-apply-Leray})
implies that $H^1(Z, f'_*\mathcal{I}') = 0$.
Since $i : Z \to X$ is affine we conclude that
$R^1i_*f'_*\mathcal{I}' = 0$ hence $H^1(X, i_*f'_*\mathcal{I}') = 0$
by Leray again. In other words, we have $H^1(X, i_*\mathcal{G}') = 0$
as desired.
\end{proof}

\end{document}
